\section{Summary}

The project has demonstrated that additive manufacturing can be effectively and economically integrated into the design of electronic device housings. The resulting prototype is a compact, self-contained device with a functional user interface, enclosed in a purpose-built housing that satisfies both mechanical and aesthetic requirements. The project also validated the practical applicability of DFMA principles in small-scale, low-volume product development, where ease of assembly, maintainability, and manufacturability are equally important to functional performance.

% The development followed a four-phase pipeline. In the system analysis phase, the board's physical dimensions and electrical architecture were thoroughly characterised to provide the foundation for all subsequent design decisions. In the mechanical design phase, the custom enclosure was modelled in Autodesk Fusion 360 as a three-component assembly, consisting of a backshell, a snap-fit front cover, and a reset button cap. DFMA principles guided the design, minimising part count and fastener requirements while optimising part orientation to reduce support material during FDM printing. The enclosure was then sliced using Ultimaker Cura and fabricated on a Creality Ender 3 V2 printer using PLA filament. In the electrical integration phase, a 3.7V LiPo battery was wired to the board's battery connector, completing the self-contained power system. Finally, in the programming phase, firmware was developed using the Arduino framework, implementing a suite of applications that include some classic games (e.g, Tetris, Alien Attack) and utility programs such as calculator and real-time clock display.

